\scanpage{53}

PS\quad Dans mon séminaire sur Riemann-Roch cette année, on doit mani\deleted{f}puler
systématiquement les immersions régulières dans le cas non noethérien,
et il apparait (hélas) que la notion introduite so\uline{u}s ce nom dans EGA
16.9.2. n'est pas la bonne notion. Un premier canular important, c'est
que si $A$ est un anneau local, et $x,y$ deux éléments de son idéal maximal,
il est facile de donner un exemple ou $(x,y)$ est une suite régulière,
et $(y,x)$ ne l'est pas. La notion de suite régulière de générateurs d'un
idéal ne jouit donc pas du tout des propriétés habituelles, comme
celle qu'il suffise de prendre une base de $J/J^2$ et de relever . Par suite,
la notion d'immersion régulière n'a pas l'air de \unsure{devoir s e|devoirs e|devoir se} comporter
comme il faut par localisation plate ou même étale, ce qui est un autre
grave inconvénient, en entrainant toute sorte d'autres d'ailleurs. Il
apparait maintenant que la notion vraiment serviable pour une suite
$(f_1,..,f_n)$ d'éléments d'un anneau (engendrant un idéal $J$) c'est que
le complexe de l'algèbre extérieure construit avec soit nul en degrés
$0$. Alors on prouve que la suite est aussi quasi-régulière, et de
plus que toute autre suite de \inserted{$n$} générateurs de $J$ a la même propriété
que celle énoncée pour $(f_1,..,f_n)$. Il y aurait lieu alors de dire que
$J$ est un idéal régulier si localement il peut être engendré par une
suite comme dessus (qu'on pourrait peut-être app\deleted{\unsure{e}}ler ``suite Koszul-régulière'' ou ``essentiellement régulière'', si on tientà ne pas changer la
terminologie des \uline{suites} régulières), et à définir de façon correspondante
la notion d'immersion régulière. Ces notions (suite K-régulière,
idéal régulier, immersion régulière) se descendent bien par fpqc. Un
composé de deux immersions régulières est itou, et les deux autres
propriétés de transitivité de 19.1.5. (iii) revu semblent marcher
encore sans hypothèse noethérienne (Berthelot en a déjà vérifié une,
% archive sheet number: 52
